\section*{Hoofdstuk \ref{ch:b2:magma-emplacement} --- Magma's en hun intrusie}

\begin{solution}{exo:b2:magma-emplacement:1}
$T = \qty{1073}{K}$. Voor $w = 1$ is $\ln w = 0$: $\log_{10}\eta = -3.545 +
9601/(1073 - 195.7) = -3.545 + 10.94 = 7.40$. Voor $w = 4$ is $\ln 4 = 1.386$:
$A = -3.545 + 1.155 = -2.390$, $B = 9601 - 3283 = 6318$, $C = 195.7 + 44.7 =
240.4$, en $\log_{10}\eta = -2.390 + 6318/832.6 = -2.390 + 7.588 = 5.20$. De
viscositeit daalt met een factor $10^{7.40 - 5.20} = 10^{2.20} \approx 160$.
\end{solution}

\begin{solution}{exo:b2:magma-emplacement:2}
Rhyoliet: $p_s = (6/0.44)^2 = \qty{186}{MPa}$, op een diepte van $\num{186e6}/(2600\times9.81)
= \qty{7.3}{km}$. Basalt: $p_s = (0.5/0.34)^2 = \qty{2.16}{MPa}$, op
$\num{2.16e6}/(2600\times9.81) = \qty{85}{m}$. Het rhyoliet begint diep in de korst bellen
te vormen, het basalt pas dicht bij het oppervlak (als het geen CO$_2$ bevat).
\end{solution}

\begin{solution}{exo:b2:magma-emplacement:3}
$(1150 + 50)/2 = \qty{600}{\celsius}$ en $(1150 + 400)/2 = \qty{775}{\celsius}$. De
\omterm{def:b2:magma-emplacement:bodies}{gang} in het koude gesteente wordt aan zijn contact sterker afgeschrikt, en krijgt de bredere
\omterm{def:b2:magma-emplacement:aureole}{afschrikrand}.
\end{solution}

\begin{solution}{exo:b2:magma-emplacement:4}
(a) Een \omterm{def:b2:magma-emplacement:bodies}{sill}, evenwijdig aan de gelaagdheid. (b) Een \omterm{def:b2:magma-emplacement:bodies}{gang}, die de gelaagdheid doorsnijdt.
(c) Een \omterm{def:b2:magma-emplacement:bodies}{batholiet}, meer dan \qty{100}{km^2} en opgebouwd uit veel intrusies.
(d) Een \omterm{def:b2:magma-emplacement:bodies}{laccoliet}, die zijn dak heeft opgebold. (e) Een \omterm{def:b2:magma-emplacement:bodies}{gang}: graniet heeft geen gelaagdheid,
dus elke plaat erin is een \omterm{def:b2:magma-emplacement:bodies}{gang}, wat haar oriëntatie ook is.
\end{solution}

\begin{solution}{exo:b2:magma-emplacement:5}
$(1 - 1.35\phi)^{-2.5}$: voor $\phi = 0.05$ is $0.9325^{-2.5} = 1.19$, dus
\qty{119}{Pa.s}, tegen $1 + 2.5\times0.05 = 1.125$, \qty{113}{Pa.s}, met de formule van
Einstein (dichtbij, zoals het hoort in de verdunde limiet). Voor $0.20$: $0.73^{-2.5} =
2.20$, \qty{220}{Pa.s}; voor $0.40$: $0.46^{-2.5} = 6.97$, \qty{697}{Pa.s}; voor $0.60$:
$0.19^{-2.5} = 63.6$, ongeveer \qty{6360}{Pa.s}. Vastlopen bij $\phi = 1/1.35 = 0.74$.
\end{solution}

\begin{solution}{exo:b2:magma-emplacement:6}
$G = 300\times9.81 = \qty{2943}{Pa/m}$. (a) $\bar u = 2943\times1^2/(12\times100) =
\qty{2.45}{m/s}$; $q = \qty{2.45}{m^2/s}$; $Re = 2700\times2.45\times1/100 = 66$, laminair;
$t = 10000/2.45 = \qty{4080}{s}$, ongeveer \qty{1.1}{h}. (b) $\bar u =
2943/(12\times\num{e6}) = \qty{2.45e-4}{m/s}$, $t = \qty{4.08e7}{s} = 472$ dagen; met
$w = \qty{10}{m}$ is $\bar u$ honderd keer groter, \qty{0.0245}{m/s}, en $t =
\qty{4.08e5}{s} = 4.7$ dagen. (c) De \omterm{def:b2:magma-emplacement:bodies}{gang} van \qty{1}{m} koelt af in ongeveer
$0.5^2/\num{e-6} = \qty{2.5e5}{s}$, 2.9 dagen; de \omterm{def:b2:magma-emplacement:bodies}{gang} van \qty{10}{m} in $5^2/\num{e-6} =
\qty{2.5e7}{s}$, 289 dagen. Het basalt (\qty{1.1}{h}) en het rhyoliet in de brede
\omterm{def:b2:magma-emplacement:bodies}{gang} (4.7 dagen) komen aan voordat hun \omterm{def:b2:magma-emplacement:bodies}{gangen} stollen; het rhyoliet in de
\omterm{def:b2:magma-emplacement:bodies}{gang} van \qty{1}{m} (472 dagen) niet.
\end{solution}

\begin{solution}{exo:b2:magma-emplacement:7}
Per kilogram \omterm{def:b2:magma-emplacement:magma}{magma} is $n = 10/18.015 = \qty{0.555}{mol}$ water en $nRT =
0.555\times8.314\times1173 = \qty{5414}{J}$; de resterende smelt neemt $0.99/2400 =
\qty{4.13e-4}{m^3}$ in. Bij \qty{100}{MPa}: $V = \qty{5.41e-5}{m^3}$, fractie
$5.41/(5.41 + 41.3) = 0.12$. Bij \qty{10}{MPa}: $V = \qty{5.41e-4}{m^3}$, fractie
$0.57$. Bij \qty{0.1}{MPa}: $V = \qty{0.0541}{m^3}$, fractie $0.992$. De waarde bij
\qty{100}{MPa} is het minst betrouwbaar: waterdamp bij zo'n druk is dicht en
verre van ideaal (en veel van het water zou in feite nog opgelost zijn).
\end{solution}

\begin{solution}{exo:b2:magma-emplacement:8}
(a) $2300 + 60z = 2650$ geeft $z = \qty{5.8}{km}$. (b) De smelt heeft
\qty{2580}{kg/m^3}: $z = 280/60 = \qty{4.7}{km}$, ongeveer \qty{1.2}{km} hoger. (c)
Bij $z = \qty{7.8}{km}$ heeft het gesteente $2300 + 60\times7.8 = \qty{2770}{kg/m^3}$; de
kracht is $(2770 - 2650)\times9.81 = \qty{1.2}{kN}$ per kubieke meter, omhoog.
\end{solution}

\begin{solution}{exo:b2:magma-emplacement:9}
$v = 2\times300\times9.81\times(1000)^2/(9\times\num{e19}) = \qty{6.5e-11}{m/s}$, of
$\num{6.5e-11}\times\num{3.156e7} = \qty{2.1e-3}{m}$ per jaar, \qty{2.1}{mm/yr}. Een stijging van
\qty{10}{km} duurt $\num{e4}/\num{2.1e-3} = \qty{4.8}{Myr}$. De tijdschaal van de afkoeling
is $(1000)^2/\num{e-6} = \qty{e12}{s} = \qty{32}{kyr}$: het lichaam stolt na
een stijging van ongeveer $32000\times0.0021 \approx \qty{70}{m}$. Een magmatische \omterm{def:b2:magma-emplacement:ascent}{diapier} kan
zo'n korst niet doorkruisen; hij heeft een hetere, zachtere omgeving nodig.
\end{solution}

\begin{solution}{exo:b2:magma-emplacement:10}
(a) $x_{\mathrm{OH}} = \sqrt{K_2x_m} = \sqrt{K_1K_2f}$, dus $x = K_1f + \frac12\sqrt{K_1K_2f}$.
De verhouding van de eerste term tot de tweede is $2\sqrt{K_1f/K_2}$: ze gaat naar nul
als $f \to 0$ (de vierkantswortel overheerst) en groeit onbegrensd met $f$ (de
lineaire term overheerst). (b) $K_1 f = \frac12\sqrt{K_1K_2 f}$ geeft $K_1^2f^2 =
K_1K_2f/4$, dus $f = K_2/(4K_1)$. (c) $f = 0.2/(\num{4e-4}) = \qty{500}{MPa}$. Bij
\qty{100}{MPa}: $x_m = 0.0100$, $\frac12\sqrt{\num{e-4}\times0.2\times100} =
\frac12\times0.0447 = 0.0224$, dus $x = 0.0324$ en het aandeel hydroxyl is
$0.0224/0.0324 = \qty{69}{\percent}$.
\end{solution}

\begin{solution}{exo:b2:magma-emplacement:11}
(a) $1000/\num{4e6} = \qty{2.5e-4}{km^3}$ per jaar, dat is $\qty{2.5e5}{m^3}$ per jaar,
of $\num{2.5e5}/\num{3.156e7} = \qty{0.0079}{m^3/s}$, ongeveer acht liter per seconde.
(b) De \omterm{def:b2:magma-emplacement:bodies}{gang} voert $1000\times2\times0.8 = \qty{1600}{m^3/s}$ aan, en levert
$\qty{e12}{m^3}$ in \qty{6.25e8}{s}, ongeveer 20 jaar: $\num{5e-6}$ van de
\qty{4}{Myr}. (c) Elke plaat bevat $20\times\num{e4}\times\num{e4} = \qty{2e9}{m^3} =
\qty{2}{km^3}$, dus 500 platen, gemiddeld één om de $\num{4e6}/500 = 8000$ jaar.
(d) $t = (10)^2/(4\times0.65^2\times\num{0.7e-6}) = \qty{8.5e7}{s}$, 2.7 jaar. Elke
plaat stolt duizenden jaren voordat de volgende aankomt: tenzij de aanvoer
in snelle stoten komt, is het \omterm{def:b2:magma-emplacement:bodies}{pluton} nooit één groot magmalichaam, zoals de
datering van grote \omterm{def:b2:magma-emplacement:bodies}{plutonen} doet vermoeden.
\end{solution}

\begin{solution}{exo:b2:magma-emplacement:12}
(a) Bij vaste $x$ is $\ln(T - T_0) = \text{const} - \frac12\ln t - x^2/(4\kappa t)$, waarvan
de afgeleide $-1/(2t) + x^2/(4\kappa t^2)$ nul wordt bij $t^* = x^2/(2\kappa)$. Dan is
$e^{-x^2/(4\kappa t^*)} = e^{-1/2}$ en $\sqrt{4\pi\kappa t^*} = x\sqrt{2\pi}$, dus $\Delta
T_{\max} = (T_m - T_0)\,b\,e^{-1/2}/(x\sqrt{2\pi}) = (T_m - T_0)\,b/(x\sqrt{2\pi e})$. (b)
$\sqrt{2\pi e} = 4.13$: $x = 1100\times2/(100\times4.13) = \qty{5.3}{m}$, bereikt na
$t^* = 5.3^2/(2\times\num{0.7e-6}) = \qty{2.0e7}{s}$, ongeveer 230 dagen (op slechts 2.7
gangbreedten is de benadering met een vlakke bron grof). (c) De verwarming neemt af
als $b/x$: een dunne \omterm{def:b2:magma-emplacement:bodies}{gang} draagt weinig warmte, en het gesteente wordt alleen
binnen enkele gangbreedten merkbaar opgewarmd.
\end{solution}

\begin{solution}{pb:b2:magma-emplacement:1}
\textbf{1.} $\log_{10}\eta = -4.55 + 5510/(1423 - 611) = -4.55 + 6.786 = 2.24$;
$\eta = \qty{172}{Pa.s}$.

\textbf{2.} $-4.55 + 5510/762 = -4.55 + 7.231 = 2.68$, $\eta = \qty{480}{Pa.s}$: een
factor $480/172 = 2.8$ voor \qty{50}{K}.

\textbf{3.} $172\times(1 - 1.35\times0.20)^{-2.5} = 172\times0.73^{-2.5} = 172\times2.20 =
\qty{378}{Pa.s}$.

\textbf{4.} $\phi = 1/1.35 = 0.74$. De veldgegevens plaatsen het einde van de stroming in bulk
bij een smeltfractie van \qtyrange{30}{50}{\percent}, dat is
\qtyrange{50}{70}{\percent} kristallen: eerder, omdat echte kristallen geen
bollen van één grootte zijn.

\textbf{5.} Basement: $2800 - 2710 = \qty{90}{kg/m^3}$, het \omterm{def:b2:magma-emplacement:magma}{magma} is lichter.
Sedimenten: $2450 - 2710 = \qty{-260}{kg/m^3}$, het \omterm{def:b2:magma-emplacement:magma}{magma} is dichter.

\textbf{6.} $p = 2450\times9.81\times3000 = \qty{72.1}{MPa}$.

\textbf{7.} $90\times9.81 = \qty{883}{Pa/m}$.

\textbf{8.} $\bar u = 883\times2.0^2/(12\times378) = 3532/4536 = \qty{0.78}{m/s}$.

\textbf{9.} $q = 0.78\times2.0 = \qty{1.56}{m^2/s}$; over \qty{1.0}{km} ongeveer
\qty{1560}{m^3/s}.

\textbf{10.} $Re = 2710\times0.78\times2.0/378 = 11$: laminair.

\textbf{11.} $12000/0.78 = \qty{15400}{s}$, ongeveer \qty{4.3}{h}.

\textbf{12.} $(1.0)^2/\num{0.7e-6} = \qty{1.43e6}{s}$, 16.5 dagen, veel langer dan
\qty{4.3}{h}: de \omterm{def:b2:magma-emplacement:bodies}{gang} blijft open zolang het \omterm{def:b2:magma-emplacement:magma}{magma} stroomt.

\textbf{13.} Op \qty{3.0}{km} ontmoet het \omterm{def:b2:magma-emplacement:magma}{magma} sedimenten die lichter zijn dan het \omterm{def:b2:magma-emplacement:magma}{magma} zelf: de
basis van de sedimenten is zijn \omterm{def:b2:magma-emplacement:ascent}{niveau van neutraal drijfvermogen}. Daarboven zou het \omterm{def:b2:magma-emplacement:magma}{magma}
dichter zijn dan zijn omgeving, dus spreidt het zich zijwaarts langs het contact uit als een
\omterm{def:b2:magma-emplacement:bodies}{sill}.

\textbf{14.} $p_s = (0.40/0.34)^2 = \qty{1.38}{MPa}$, bereikt op
$\num{1.38e6}/(2450\times9.81) = \qty{58}{m}$ onder het oppervlak.

\textbf{15.} $0.34\times\sqrt{72.1} = 0.34\times8.49 = \qty{2.9}{\percent}$, veel meer dan
de aanwezige \qty{0.40}{\percent}: het \omterm{def:b2:magma-emplacement:magma}{magma} van de \omterm{def:b2:magma-emplacement:bodies}{sill} ontmengt geen water.

\textbf{16.} CO$_2$ is veel minder oplosbaar dan water: het raakt verzadigd en vormt al diep
tijdens de opstijging bellen, dus het eerste gas van een basalt is vooral CO$_2$.

\textbf{17.} Per kilogram is $n = 4.0/18.015 = \qty{0.222}{mol}$, $V = 0.222\times8.314\times1423/
\num{e5} = \qty{0.0263}{m^3}$; de smelt neemt $0.996/2710 = \qty{3.68e-4}{m^3}$ in; de
fractie is $0.0263/(0.0263 + 0.000368) = 0.986$, \qty{98.6}{\percent}.

\textbf{18.} $(1150 + 100)/2 = \qty{625}{\celsius}$.

\textbf{19.} Na 10 jaar is $2\sqrt{\kappa t} = 2\sqrt{\num{0.7e-6}\times\num{3.156e8}} =
\qty{29.7}{m}$. Het punt ligt op $x = \qty{60}{m}$ van het midden: $T = 100 +
525\,[\operatorname{erf}(-10/29.7) + \operatorname{erf}(110/29.7)] = 100 + 525\,(-0.365 +
1.000) = \qty{433}{\celsius}$.

\textbf{20.} In het midden is $T = 100 + 1050\operatorname{erf}(50/2\sqrt{\kappa t})$: na
10 jaar $1050\operatorname{erf}(1.68) = 1050\times0.983$, $T = \qty{1132}{\celsius}$; na
50 jaar is $2\sqrt{\kappa t} = \qty{66.5}{m}$, $1050\operatorname{erf}(0.752) = 1050\times0.713$,
$T = \qty{848}{\celsius}$.

\textbf{21.} Bij $x = \qty{130}{m}$: \qty{294}{\celsius} na 300 jaar,
\qty{295}{\celsius} na 363 jaar (het maximum), \qty{295}{\celsius} na 420
jaar. Zonder latente warmte bereikt het gesteente \qty{80}{m} boven de \omterm{def:b2:magma-emplacement:bodies}{sill} nooit helemaal
\qty{300}{\celsius}.

\textbf{22.} $\num{9e8}\times1.772/(\num{2.5e6}\times1050) = 0.608$.

\textbf{23.} $\lambda = 0.64$: $e^{-0.4096}/(0.64\times1.6346) = 0.6639/1.0461 = 0.635$;
$0.65$: $0.6554/1.0673 = 0.614$; $0.66$: $0.6469/1.0886 = 0.594$. De waarde $0.608$
ligt tussen $0.65$ en $0.66$, dicht bij $0.65$: $\lambda = 0.65$.

\textbf{24.} $100 + 1050/(1 + 0.642) = 100 + 640 = \qty{740}{\celsius}$. De latente
warmte die aan het stollingsfront vrijkomt, houdt de kant van het \omterm{def:b2:magma-emplacement:magma}{magma} langer op
\qty{1150}{\celsius}, dus het contact ontvangt meer warmte.

\textbf{25.} $(50)^2/\num{0.7e-6} = \qty{3.57e9}{s}$, 113 jaar.

\textbf{26.} $t = (100)^2/(16\times0.65^2\times\num{0.7e-6}) = \num{e4}/\num{4.73e-6} =
\qty{2.11e9}{s}$, dus $\num{2.11e9}/\num{3.156e7}$: de \omterm{def:b2:magma-emplacement:bodies}{sill} is vast na
\textbf{\qty[text-series-to-math=true]{67}{} jaar}. Dat is korter dan
$(b/2)^2/\kappa$ met de factor $1/(4\lambda^2) = 0.59$: die tijdschaal geeft alleen
de grootteorde, en beide wanden van de \omterm{def:b2:magma-emplacement:bodies}{sill} voeren tegelijk warmte af naar een koud gastgesteente,
snel genoeg om ook de latente warmte af te voeren.
\end{solution}
